\section{SCOPE}\label{sec:method}

\subsection{Design Rationale}

\hyre{} mirrors a small piece of how an expert reads a legal scenario. A
practitioner facing a fact pattern forms a quick intuition about what is at
issue (e.g. proximate cause, a hearsay exception, or a state-specific notice rule),
looks up the controlling authority by that intuition, and then judges the
evidence without allowing intuition to bias the judgment.
\hyre{} encodes the same loop. The first call asks the model to form an
initial draft answer from its priors and to write a separate legal search passage
for retrieval. The second call reasons only over the question and the
retrieved evidence; it never sees the draft answer, so the final answer is judged
against what the corpus actually says.
Figure~\ref{fig:pipeline_art} shows the resulting two-call pipeline.

\subsection{Two Calls, One Retrieval Query}\label{sec:method:snaphyre}

\hyre{} separates query formation from final answering. The same model
$G_\theta$ is used for both calls; only the prompt template differs (the
superscript). Given a formatted legal question $q$, the first call produces
a private draft answer $a_0$ and a legal search passage $p$:

\begin{equation}
(a_0, p) = G_{\theta}^{\mathrm{snap}}(q).
\end{equation}

Only $p$ is exposed to retrieval. Let
$\mathrm{enc} : \mathcal{T} \to \mathbb{R}^d$ denote the dense encoder shared
by queries and corpus passages, and let $\mathcal{C}$ denote the searchable
corpus. The retriever ranks every $c \in \mathcal{C}$ by similarity to the
embedded search passage and returns the top-$k$ candidates,

\begin{equation}
E_k \;=\; \operatorname*{TopK}_{c \in \mathcal{C}} \;\bigl\langle \mathrm{enc}(p),\, \mathrm{enc}(c) \bigr\rangle,
\end{equation}

after which a cross-encoder reranker reorders $E_k$ on the same passage $p$.
The final answer is then

\begin{equation}
\hat{a} = G_{\theta}^{\mathrm{answer}}(q, E_k),
\end{equation}

with $a_0$ withheld from the final context. The first call can use the model's
prior knowledge of law to choose vocabulary; the second call must answer from
the retrieved corpus text. Algorithm~\ref{alg:scope} summarizes the pipeline.

\begin{algorithm}[t]
\caption{\hyre{} pipeline for one legal question.}\label{alg:scope}
\begin{algorithmic}[1]
\REQUIRE question $q$; generator $G_\theta$; encoder $\mathrm{enc}$; corpus $\mathcal{C}$; depth $k$
\STATE $(a_0,\, p) \gets G_{\theta}^{\mathrm{snap}}(q)$ \hfill // first call
\STATE \textbf{validate} $a_0$ has a parseable answer label and $p$ contains no answer marker
\STATE $E_k \gets \operatorname*{TopK}_{c \in \mathcal{C}} \langle \mathrm{enc}(p), \mathrm{enc}(c) \rangle$ \hfill // retrieve using $p$ only
\STATE \textbf{rerank} $E_k$ by cross-encoder on $p$
\STATE $\hat{a} \gets G_{\theta}^{\mathrm{answer}}(q, E_k)$ \hfill // second call; $a_0$ is discarded
\STATE \textbf{return} $\hat{a}$
\end{algorithmic}
\end{algorithm}

\paragraph{Contrast with HyDE.}
HyDE writes a single hypothetical document directly from the question,
\begin{equation}
p_{\mathrm{hyde}} = G_{\theta}^{\mathrm{hyde}}(q), \qquad
E_k^{\mathrm{hyde}} = R_k(p_{\mathrm{hyde}}),
\end{equation}
and uses $p_{\mathrm{hyde}}$ as the retrieval handle \citep{gao2023hyde}. \hyre{}
keeps that generated-query premise but reshapes the first call: the model
first commits to a tentative legal frame $a_0$, then writes $p$ as language
consistent with that frame. The snap draft acts as an information bottleneck
that selects \emph{what kind} of authority the passage should look like, but
it is discarded before the corpus is queried and never reaches the answer
call. For a torts fact pattern about a broken leg, crutches, and a later
fall, this turns an accident-heavy query into a passage about proximate cause
and foreseeable subsequent injury
(Figure~\ref{fig:query_contrast_diagram}).

\begin{figure*}[t]
  \centering
  \includegraphics[width=\textwidth]{22_query_contrast_diagram.png}
  \caption{What the two methods send to the retriever for one BarExamQA item
  (\texttt{mbe\_1175}). Raw-question RAG embeds the narrative
  fact pattern; SCOPE embeds a short doctrinal passage about foreseeable
  intervening injury, which is the language of the rule the retriever
  actually needs to surface.}
  \label{fig:query_contrast_diagram}
\end{figure*}

\subsection{Grounding Guardrails}

The two-call structure gives the method two guardrails that follow from the
formulation above. First, the draft $a_0$ is private: it can shape what to
look for, but it is not handed in as evidence. The second call's answer
depends on the draft only through the search passage and the retrieved set,

\begin{equation}
\hat{a} \;\perp\!\!\!\perp\; a_0 \;\bigm|\; (q, E_k),
\end{equation}

so the model cannot anchor its evidence-stage reasoning on a guess it made
without seeing the corpus. Second, the search passage $p$ is written like
legal authority rather than chain-of-thought reasoning: it names the rule,
statutory condition, or doctrinal test that would resolve the question, not
the steps the model used to get there. The intuition step is allowed to use
model priors and the evidence step is forced back to the corpus.

Each \hyre{} response follows three steps: validate that the first call
produced both a formatted draft answer and a clean search passage; retrieve
and rerank using only the search passage; answer the original question from
the retrieved passages without exposing the draft answer. The retrieval
passage is constrained to reflect the corpus it is searching over,
typically describing the controlling rule, statutory condition, legal
relationship, or doctrinal test without answer labels or an
\texttt{Answer:} header, as to prevent answer leakage from the first step to the second step.